% ===========================================================================
%  Raport — Studiu Preliminar (Faza 0)
%  Antrenare și Cuantizare FP8 pe ViT-Tiny / CIFAR-10
%  Tudose Alexandru · Lucrare de Licență 2026
% ===========================================================================
\documentclass[12pt,a4paper]{article}

% --- Encoding & fonts (XeTeX / tectonic) ---
\usepackage{fontspec}
\defaultfontfeatures{Ligatures=TeX}

% --- Language ---
\usepackage{polyglossia}
\setmainlanguage{romanian}

% --- Page geometry ---
\usepackage[
  top=2.5cm, bottom=2.5cm,
  left=2.5cm, right=2.5cm
]{geometry}

% --- Math ---
\usepackage{amsmath, amssymb}

% --- Graphics ---
\usepackage{graphicx}
\usepackage{float}
\usepackage{caption}
\usepackage{subcaption}
\captionsetup{font=small, labelfont=bf}

% --- Tables ---
\usepackage{booktabs}
\usepackage{array}
\usepackage{multirow}

% --- Colors ---
\usepackage{xcolor}
\definecolor{fp32color}{HTML}{0072B2}
\definecolor{fp16color}{HTML}{009E73}
\definecolor{fp8color}{HTML}{E69F00}
\definecolor{lightgray}{HTML}{F4F4F4}

% --- Hyperlinks ---
\usepackage[
  colorlinks=true,
  linkcolor=fp32color,
  citecolor=fp32color,
  urlcolor=fp32color
]{hyperref}

% --- Misc ---
\usepackage{microtype}
\usepackage{parskip}
\setlength{\parindent}{0pt}
\setlength{\parskip}{6pt}
\usepackage{enumitem}

% --- Header / Footer ---
\usepackage{fancyhdr}
\pagestyle{fancy}
\fancyhf{}
\lhead{\small Tudose Alexandru --- Studiu Preliminar (Faza 0)}
\rhead{\small 2026}
\cfoot{\thepage}
\renewcommand{\headrulewidth}{0.4pt}

% ===========================================================================
\begin{document}
% ===========================================================================

% --- Title page ---
\begin{titlepage}
  \centering
  {\large Universitatea Transilvania din Brașov\\}
  {\small Facultatea de Matematică și Informatică\\}
  {\small Specializarea: Informatică Aplicată\\[2cm]}

  \rule{\linewidth}{0.4pt}\\[0.5cm]
  {\LARGE\bfseries Studiu Preliminar (Faza 0)\\[0.3cm]}
  {\Large\bfseries Antrenare ViT-Tiny pe CIFAR-10\\[0.2cm]}
  {\Large\bfseries Cuantizare FP16, FP8 E4M3FN\\[0.2cm]}
  {\large Proof of Concept}\\[0.3cm]
  \rule{\linewidth}{0.4pt}\\[2cm]

  \begin{minipage}{0.45\textwidth}
    \textbf{Student:}\\
    Tudose Alexandru\\
    Grupa 10LF333\\
    An III, Informatică Aplicată
  \end{minipage}
  \hfill
  \begin{minipage}{0.45\textwidth}
    \raggedleft
    \textbf{Coordonator:}\\
    Conf.\ dr.\ ing.\ Honorius Gâlmeanu\\
    Departamentul de Informatică
  \end{minipage}\\[3cm]

  {\small Brașov, 2026}
\end{titlepage}

% --- TOC ---
\tableofcontents
\newpage

% ===========================================================================
\section{Introducere}
% ===========================================================================

Acest raport documentează \textbf{studiul preliminar} (Faza~0) al lucrării de
licență, care a servit drept \emph{proof of concept} pentru viabilitatea
cuantizării numerice pe Vision Transformers. Experimentele au fost concepute
pentru a valida pipeline-ul experimental și a demonstra fezabilitatea cuantizării
înainte de trecerea la evaluarea pe ImageNet cu model pretrained.

\textbf{Obiective:}
\begin{enumerate}
  \item Antrenarea unui ViT-Tiny pe CIFAR-10 cu diferite configurații de
    precizie (FP32 vs.\ FP16) și augmentare (basic vs.\ extended).
  \item Evaluarea impactului fiecărui factor asupra acurateții, generalizării
    și timpului de antrenare.
  \item Testarea cuantizării post-training FP8 (\texttt{float8\_e4m3fn}) pe
    cel mai bun model antrenat.
\end{enumerate}

\textbf{Statut:} Studiul preliminar este \textbf{finalizat}. Concluziile sale
au informat direcția fazelor ulterioare (1--3) ale lucrării. A fost identificat
faptul că abordarea FP8 nativă (\texttt{torch.float8\_e4m3fn}) este
problematică, ceea ce a condus la adoptarea cuantizării INT8 cu scalare liniară
în fazele principale.

\textbf{Notă:} Acest raport este inclus în lucrarea de licență ca \emph{studiu
preliminar / test funcțional}. Rezultatele principale ale lucrării provin din
Fazele 1--3 (evaluare pe ImageNet/ImageNette cu model pretrained).

% ===========================================================================
\section{Arhitectura Modelului și Dataset}
% ===========================================================================

\subsection{Vision Transformer --- ViT-Tiny}

Modelul utilizat este \texttt{vit\_tiny\_patch16\_224} din biblioteca
\texttt{timm}~\cite{timm}, cu 5.7 milioane de parametri.

\begin{table}[H]
  \centering
  \caption{Parametrii arhitecturali ai ViT-Tiny}
  \begin{tabular}{ll}
    \toprule
    \textbf{Proprietate} & \textbf{Valoare} \\
    \midrule
    Parametri totali    & 5.7 milioane \\
    Patch size          & $16 \times 16$ pixeli \\
    Input resolution    & $224 \times 224$ \\
    Embedding dim       & 192 \\
    Număr blocuri       & 12 \\
    Număr capete (MHA)  & 3 \\
    MLP ratio           & 4 (hidden dim = 768) \\
    \bottomrule
  \end{tabular}
  \label{tab:vit_config_prelim}
\end{table}

\subsection{CIFAR-10}

CIFAR-10 conține 60\,000 imagini de $32 \times 32$ pixeli, distribuite în
10 clase. Imaginile sunt redimensionate la $224 \times 224$ pentru
compatibilitate cu ViT. Setul este împărțit în 50\,000 de imagini de antrenare
și 10\,000 de imagini de validare.

\textbf{Limitare:} CIFAR-10 este un dataset simplu și mic, insuficient pentru
a trage concluzii despre comportamentul cuantizării pe sarcini complexe.
Rezultatele au fost utilizate exclusiv ca validare a pipeline-ului experimental.

% ===========================================================================
\section{Design Experimental}
% ===========================================================================

\subsection{Matricea experimentelor}

Au fost realizate 4 experimente combinând 2 factori:

\begin{table}[H]
  \centering
  \caption{Matricea experimentelor $2 \times 2$}
  \begin{tabular}{lcc}
    \toprule
    & \textbf{Basic augmentation} & \textbf{Extended augmentation} \\
    \midrule
    \textbf{FP32} & BaseFP32 & AugmFP32 \\
    \textbf{FP16 (AMP)} & BaseFP16 & AugmFP16 \\
    \bottomrule
  \end{tabular}
  \label{tab:experiment_matrix}
\end{table}

\textbf{Augmentare basic:} Random horizontal flip + normalizare standard.

\textbf{Augmentare extended:} Random crop, horizontal flip, color jitter,
random erasing, CutMix/MixUp, label smoothing.

\subsection{Configurație antrenare}

\begin{table}[H]
  \centering
  \caption{Hiperparametri comuni}
  \begin{tabular}{ll}
    \toprule
    \textbf{Parametru} & \textbf{Valoare} \\
    \midrule
    Epoci               & 50 \\
    Batch size           & 64 \\
    Optimizer            & AdamW \\
    Learning rate        & $3 \times 10^{-4}$ \\
    Scheduler            & Cosine annealing (warm restarts) \\
    Pre-training         & ImageNet-1k (\texttt{pretrained=True}) \\
    Fine-tuning dataset  & CIFAR-10 \\
    Device               & Apple M4 MPS \\
    PyTorch              & 2.9.1 \\
    \bottomrule
  \end{tabular}
  \label{tab:training_config}
\end{table}

\subsection{Hardware}

Toate experimentele au fost executate pe MacBook Air cu Apple M4, 16~GB RAM,
utilizând backend-ul MPS (Metal Performance Shaders) din PyTorch.

% ===========================================================================
\section{Rezultate Antrenare}
% ===========================================================================

\subsection{Comparație globală}

\begin{table}[H]
  \centering
  \caption{Rezultate finale --- toate experimentele}
  \begin{tabular}{lrrrrr}
    \toprule
    \textbf{Experiment} & \textbf{Best Val Acc} & \textbf{Final Train Acc} &
    \textbf{Overfit Gap} & \textbf{Conv.\ Epoch} & \textbf{Timp (h)} \\
    \midrule
    BaseFP32 & 78.99\% & 97.29\% & 18.37 pp & 43 & 7.26 \\
    BaseFP16 & 79.19\% & 97.43\% & 18.40 pp & 35 & 7.45 \\
    AugmFP32 & 82.88\% & 86.80\% & \phantom{0}4.09 pp & 45 & 7.38 \\
    AugmFP16 & \textbf{83.41\%} & 85.67\% & \textbf{\phantom{0}2.26 pp} & 48 & \textbf{7.09} \\
    \bottomrule
  \end{tabular}
  \label{tab:all_results}
\end{table}

\textit{Overfit Gap = Final Train Acc $-$ Final Val Acc. Conv.\ Epoch = prima
epocă ce atinge 99.5\% din best val acc.}

\begin{figure}[H]
  \centering
  \includegraphics[width=\textwidth]{../results/preliminary/report_plots/A1_val_acc_all.png}
  \caption{Curbele de acuratețe de validare pentru cele 4 experimente pe parcursul
    a 50 de epoci. Configurațiile cu augmentare extinsă (Augm*) ating o acuratețe
    semnificativ mai mare decât cele cu augmentare basic (Base*). Diferența între
    FP32 și FP16 este neglijabilă la același tip de augmentare.}
  \label{fig:val_acc_all}
\end{figure}

\begin{figure}[H]
  \centering
  \includegraphics[width=\textwidth]{../results/preliminary/report_plots/A2_train_val_acc_grid.png}
  \caption{Grid $2 \times 2$: acuratețe training vs.\ validare pentru fiecare
    experiment. Zona umbrită reprezintă gap-ul de overfitting. Experimentele
    Base* prezintă un gap masiv ($\approx 18$~pp), în timp ce Augm* reduc drastic
    overfitting-ul la 2--4~pp.}
  \label{fig:train_val_grid}
\end{figure}

\subsection{Efectul augmentării}

Augmentarea extinsă este factorul dominant:

\begin{itemize}
  \item \textbf{FP32:} Basic $\to$ Extended: $+3.90$~pp acuratețe
  \item \textbf{FP16:} Basic $\to$ Extended: $+4.21$~pp acuratețe
  \item Overfitting gap redus de la $\approx 18.4$~pp la $\approx 2.3$--4.1~pp
\end{itemize}

\begin{figure}[H]
  \centering
  \begin{subfigure}[b]{0.48\textwidth}
    \includegraphics[width=\textwidth]{../results/preliminary/report_plots/C1_best_val_acc_bar.png}
    \caption{Acuratețe de validare maximă}
  \end{subfigure}
  \hfill
  \begin{subfigure}[b]{0.48\textwidth}
    \includegraphics[width=\textwidth]{../results/preliminary/report_plots/C2_overfitting_gap_bar.png}
    \caption{Gap de overfitting}
  \end{subfigure}
  \caption{Comparație bar chart. (a) Augmentarea extinsă câștigă $\approx 4$~pp
    acuratețe. (b) Overfitting-ul scade dramatic cu augmentare extinsă, de la
    18.4~pp la 2.3~pp.}
  \label{fig:acc_overfit_bars}
\end{figure}

\subsection{Efectul preciziei numerice}

Diferența între FP32 și FP16 (AMP) este neglijabilă:
\begin{itemize}
  \item Cu augmentare basic: FP16 este marginal mai bun cu $+0.21$~pp
  \item Cu augmentare extended: FP16 este mai bun cu $+0.52$~pp
  \item Timpii de antrenare sunt comparabili ($\pm 3$\%)
\end{itemize}

\textbf{Concluzie:} Precizia numerică FP32 vs.\ FP16 nu are impact practic
asupra acurateții pe CIFAR-10, confirmând robustețea ViT la reducerea preciziei.

\subsection{Curbe de loss}

\begin{figure}[H]
  \centering
  \includegraphics[width=\textwidth]{../results/preliminary/report_plots/B_loss_curves.png}
  \caption{Curbele de loss (training și validare) pentru toate experimentele.
    Experimentele Base* converg la un loss de validare ridicat ($\approx 0.9$),
    cu un gap mare față de training loss ($\approx 0.08$), indicând overfitting sever.
    Experimentele Augm* mențin un gap mic între train și val loss.}
  \label{fig:loss_curves}
\end{figure}

\subsection{Timing}

\begin{figure}[H]
  \centering
  \includegraphics[width=0.65\textwidth]{../results/preliminary/report_plots/C3_training_time_bar.png}
  \caption{Timpul total de antrenare. Toate experimentele durează
    $\approx 7$--7.5~ore pe Apple M4 MPS. AugmFP16 este marginal cel mai rapid
    (7.09~h), deoarece FP16 reduce overhead-ul per operație.}
  \label{fig:training_time}
\end{figure}

\begin{table}[H]
  \centering
  \caption{Timing per epocă}
  \begin{tabular}{lrr}
    \toprule
    \textbf{Experiment} & \textbf{Medie (sec)} & \textbf{Std (sec)} \\
    \midrule
    BaseFP32 & 522.90 & $\pm 16.45$ \\
    BaseFP16 & 536.22 & $\pm 21.99$ \\
    AugmFP32 & 531.33 & $\pm 27.68$ \\
    AugmFP16 & 510.29 & $\pm 18.02$ \\
    \bottomrule
  \end{tabular}
  \label{tab:epoch_timing}
\end{table}

% ===========================================================================
\section{Cuantizare Post-Training FP8}
% ===========================================================================

\subsection{Metodologie}

Cel mai bun model antrenat (\textbf{AugmFP16}, acuratețe 83.41\%) a fost supus
cuantizării post-training la \texttt{float8\_e4m3fn} la două granularități:

\begin{enumerate}
  \item \textbf{Per-tensor:} Un singur factor de scalare per matrice de ponderi.
    \begin{equation}
      s = \frac{\text{FP8\_MAX}}{\max(|W|)}, \quad
      \hat{W} = \text{clamp}(W \cdot s,\, {-448},\, {448}).\text{to(fp8)}.\text{to(dtype)} / s
    \end{equation}
  \item \textbf{Per-channel:} Un factor de scalare per rând (dim~0) al matricii.
\end{enumerate}

\textbf{Notă importantă:} Această abordare (cast direct la \texttt{float8\_e4m3fn})
a fost ulterior identificată drept \emph{problematică} de către coordonator,
deoarece nu toate valorile FP32 sunt reprezentabile în formatul FP8. Fazele
ulterioare (1--3) au adoptat cuantizarea INT8 cu scalare liniară, care
garantează o mapare bijectivă la 256 valori discrete.

\subsection{Rezultate}

\begin{table}[H]
  \centering
  \caption{Rezultate cuantizare FP8 pe modelul AugmFP16}
  \begin{tabular}{lrrr}
    \toprule
    \textbf{Metodă} & \textbf{Acuratețe} & \textbf{Degradare} & \textbf{Loss} \\
    \midrule
    FP16 (original)   & 83.43\% & ---         & 0.5572 \\
    FP8 per-tensor    & 83.41\% & $+0.02$~pp  & 0.5581 \\
    FP8 per-channel   & 83.42\% & $+0.01$~pp  & 0.5577 \\
    \bottomrule
  \end{tabular}
  \label{tab:fp8_results}
\end{table}

\begin{figure}[H]
  \centering
  \includegraphics[width=0.75\textwidth]{../results/preliminary/report_plots/D_fp8_impact.png}
  \caption{Impactul cuantizării FP8 asupra acurateții. Degradarea este
    esențialmente zero ($<0.02$~pp), demonstrând viabilitatea cuantizării
    pe 8 biți pe acest model și dataset.}
  \label{fig:fp8_impact}
\end{figure}

\subsection{Discuție}

Degradarea de $+0.02$~pp (per-tensor) și $+0.01$~pp (per-channel) este
neglijabilă și se încadrează în variabilitatea numerică. Acest rezultat a
confirmat ipoteza inițială: ViT-Tiny este robust la cuantizarea pe 8 biți.

Totuși, aceste rezultate trebuie interpretate cu prudență:
\begin{itemize}
  \item CIFAR-10 este un dataset simplu --- degradarea pe sarcini mai complexe
    poate fi semnificativ mai mare.
  \item Metoda de cuantizare FP8 (\texttt{float8\_e4m3fn} direct) nu oferă
    aceeași garanție de mapare bijectivă ca INT8 cu scalare liniară.
  \item Toate ponderile au fost cuantizate, inclusiv LayerNorm și head ---
    în fazele ulterioare, acestea sunt excluse.
\end{itemize}

% ===========================================================================
\section{Monitorizare Hardware}
% ===========================================================================

\begin{table}[H]
  \centering
  \caption{Utilizare hardware în timpul antrenării}
  \begin{tabular}{lrrrr}
    \toprule
    \textbf{Experiment} & \textbf{CPU Avg} & \textbf{CPU Max} &
    \textbf{RAM Avg} & \textbf{RAM Max} \\
    \midrule
    BaseFP32 & 2.6\% & 54.3\% & 83.6\% & 89.3\% \\
    BaseFP16 & 2.6\% & 58.9\% & 81.6\% & 86.4\% \\
    AugmFP32 & 5.2\% & 59.1\% & 86.0\% & 90.8\% \\
    AugmFP16 & 4.8\% & 62.6\% & 81.6\% & 86.8\% \\
    \bottomrule
  \end{tabular}
  \label{tab:hardware}
\end{table}

\textbf{Observații:}
\begin{itemize}
  \item CPU-ul este utilizat minimal ($<5\%$ medie), indicând o sarcină
    GPU-bound pe MPS.
  \item Augmentarea extinsă crește utilizarea CPU de $\approx 2\times$
    (preprocessing suplimentar).
  \item Memoria RAM este saturată ($81\text{--}91\%$), sugerând că batch-ul de 64
    este aproape de limita celor 16~GB.
  \item \textbf{Niciun experiment nu a declanșat thermal throttling.}
\end{itemize}

\begin{figure}[H]
  \centering
  \begin{subfigure}[b]{0.48\textwidth}
    \includegraphics[width=\textwidth]{../results/preliminary/report_plots/E1_cpu_memory_bar.png}
    \caption{CPU și RAM per experiment}
  \end{subfigure}
  \hfill
  \begin{subfigure}[b]{0.48\textwidth}
    \includegraphics[width=\textwidth]{../results/preliminary/report_plots/E2_gpu_power_profile.png}
    \caption{Profilul de putere GPU/CPU}
  \end{subfigure}
  \caption{(a) Utilizare CPU și RAM medie/maximă. (b) Profilul de consum
    energetic în timp pentru AugmFP32 vs.\ AugmFP16 --- profiluri similare,
    confirmând că FP16 nu introduce overhead suplimentar pe Apple M4.}
  \label{fig:hardware}
\end{figure}

% ===========================================================================
\section{Limitări ale Studiului Preliminar}
% ===========================================================================

\begin{enumerate}
  \item \textbf{Dataset:} CIFAR-10 este prea simplu și mic (60\,000 imagini,
    10 clase) pentru a trage concluzii valide despre comportamentul cuantizării
    pe sarcini reale. Rezultatele nu se generalizează direct la ImageNet.

  \item \textbf{Metoda FP8:} Cast-ul direct la \texttt{float8\_e4m3fn} nu
    garantează o mapare bijectivă --- anumite valori FP32 nu sunt reprezentabile
    în FP8 E4M3FN. Fazele ulterioare adoptă INT8 cu scalare liniară.

  \item \textbf{Cuantizare neselectivă:} Toate ponderile au fost cuantizate,
    inclusiv LayerNorm și classification head, care sunt sensibile la erori
    numerice. Abordarea corectă (Faza 2+) exclude aceste straturi.

  \item \textbf{Fine-tuning, nu inferență pură:} Studiul preliminar antrenează
    modelul pe CIFAR-10, deci evaluează cuantizarea \emph{după fine-tuning}.
    Fazele principale evaluează cuantizarea pe un model pretrained
    \emph{fără nicio re-antrenare} (post-training quantization).

  \item \textbf{Hardware:} Experimentele pe M4 nu evaluează speedup-ul real
    al cuantizării pe hardware cu suport nativ (NVIDIA Tensor Cores).
\end{enumerate}

% ===========================================================================
\section{Concluzii și Lecții Învățate}
% ===========================================================================

\subsection{Concluzii}

\begin{enumerate}
  \item \textbf{Augmentarea domină:} Augmentarea extinsă produce un câștig de
    $+3.9$--$4.2$~pp acuratețe, reducând overfitting-ul de la 18.4~pp la 2.3~pp.
    Acesta este factorul cu cel mai mare impact asupra performanței.

  \item \textbf{FP32 vs.\ FP16 --- diferență neglijabilă:} Precizia numerică
    nu afectează semnificativ acuratețea ($\Delta < 0.5$~pp), confirmând
    robustețea ViT la \texttt{float16}.

  \item \textbf{Configurația optimă: AugmFP16} --- acuratețe 83.41\%,
    overfitting gap 2.26~pp, timp minim (7.09~h).

  \item \textbf{FP8 viabil ca proof of concept:} Degradarea $<0.02$~pp pe
    CIFAR-10, dar metoda FP8 nativă prezintă limitări teoretice.
\end{enumerate}

\subsection{Lecții pentru fazele ulterioare}

\begin{itemize}
  \item Trecerea de la CIFAR-10 la \textbf{ImageNette2-320} (proxy ImageNet)
    pentru evaluări mai relevante.
  \item Adoptarea \textbf{INT8 cu scalare liniară} în locul FP8 nativ --- mapare
    bijectivă garantată.
  \item \textbf{Cuantizare selectivă:} excluderea LayerNorm, cls\_token,
    pos\_embed și head.
  \item Evaluare în regim \textbf{post-training quantization} (fără
    re-antrenare), pe model pretrained.
  \item Analiza \textbf{per-channel vs.\ per-tensor} și \textbf{sensitivity
    analysis} per strat.
\end{itemize}

% ===========================================================================
\begin{thebibliography}{9}
% ===========================================================================

\bibitem{dosovitskiy2021vit}
A. Dosovitskiy, L. Beyer, A. Kolesnikov et al.,
``An Image is Worth 16x16 Words: Transformers for Image Recognition at Scale,''
\emph{ICLR}, 2021. \href{https://arxiv.org/abs/2010.11929}{arXiv:2010.11929}.

\bibitem{touvron2021deit}
H. Touvron, M. Cord, M. Douze et al.,
``Training data-efficient image transformers \& distillation through attention,''
\emph{ICML}, 2021. \href{https://arxiv.org/abs/2012.12877}{arXiv:2012.12877}.

\bibitem{micikevicius2022fp8}
P. Micikevicius, D. Stosic, N. Burgess et al.,
``FP8 Formats for Deep Learning,''
2022. \href{https://arxiv.org/abs/2209.05433}{arXiv:2209.05433}.

\bibitem{timm}
R. Wightman, ``PyTorch Image Models (timm),'' 2019.
\url{https://github.com/rwightman/pytorch-image-models}.

\bibitem{cifar10}
A. Krizhevsky, ``Learning Multiple Layers of Features from Tiny Images,''
Technical Report, University of Toronto, 2009.

\end{thebibliography}

\end{document}
